\documentclass[11pt,a4paper]{article}

\usepackage[a4paper,margin=1in]{geometry}
\usepackage{graphicx}
\usepackage{booktabs}
\usepackage{hyperref}
\usepackage{enumitem}
\usepackage{microtype}
\usepackage{parskip}
\usepackage{titlesec}
\usepackage{xcolor}
\usepackage{listings}

\hypersetup{
  colorlinks=true,
  linkcolor=black,
  urlcolor=blue,
  citecolor=black
}

\titleformat{\section}{\large\bfseries}{\thesection}{0.75em}{}
\titleformat{\subsection}{\normalsize\bfseries}{\thesubsection}{0.75em}{}

% --- HEADER FORMATTING ---
% Standardizing spacing: {left}{before}{after}
\titleformat{\section}{\large\bfseries}{\thesection}{0.75em}{}
\titlespacing*{\section}{0pt}{1.5ex plus 1ex minus .2ex}{1ex plus .2ex}

\titleformat{\subsection}{\normalsize\bfseries}{\thesubsection}{0.75em}{}
\titlespacing*{\subsection}{0pt}{1.2ex plus 1ex minus .2ex}{0.8ex plus .2ex}

% --- TITLE DATA ---
\title{\textbf{CodeQuery API: Technical Report}}
\author{
  \textbf{Sachin Sindhe}\\
  COMP3011 -- Web Services and Web Data\\[0.5em]
  \textbf{GitHub Repository:} \href{https://github.com/Sachin-1712/RepoLens}{Sachin-1712/RepoLens}\\
  \textbf{API Documentation:} \href{https://github.com/Sachin-1712/RepoLens/blob/main/api-docs.pdf}{View PDF Documentation}
}
\date{}

\begin{document}
\maketitle

\vspace{-0.5em}
\begin{center}
\textbf{Abstract.}
This report describes CodeQuery, a database-backed RESTful API that ingests Git repositories and supports natural-language questions about the codebase using Retrieval-Augmented Generation (RAG). The system stores repositories, analysis artefacts, and question history in PostgreSQL, performs semantic retrieval using pgvector, and optionally generates grounded answers via a local LLM (Ollama). A ``Lightweight Mode'' runs by default (API + database only), ensuring storage-safe operation on consumer hardware and graceful degradation when advanced services are unavailable.
\end{center}

\section{Introduction and Objectives}
CodeQuery is an AI-powered code intelligence API designed to ingest Git repositories and help users understand unfamiliar codebases by asking questions such as ``How is routing implemented?'' or ``Where is authentication handled?''. The system follows a Retrieval-Augmented Generation (RAG) workflow: repository files are chunked and embedded, relevant code excerpts are retrieved for a given question, and an answer is generated from that evidence (optionally using a local LLM).

The project objectives were:
\begin{itemize}[leftmargin=1.2em]
  \item Build a RESTful API meeting coursework requirements: multiple CRUD endpoints, persistent SQL storage, correct HTTP status codes, JSON responses, and OpenAPI (Swagger) documentation.
  \item Implement repository ingestion: clone a repository, discover and filter files, chunk code, and store structured chunks in the database.
  \item Provide repository Q\&A using semantic retrieval over stored chunks, with an optional local LLM for natural-language answer generation.
  \item Support \textbf{Lightweight Mode} by default to run reliably on limited storage (under 20GB) without paid APIs or heavyweight background services, while gracefully degrading when optional components are not enabled.
\end{itemize}

\section{Technology Stack and Design Rationale}

\subsection{Python and FastAPI}
Python 3.12 was selected due to its strong ecosystem for repository manipulation and semantic tooling. FastAPI was chosen because it provides:
\begin{itemize}[leftmargin=1.2em]
  \item \textbf{Asynchronous support} for I/O-heavy operations (database queries, HTTP calls).
  \item \textbf{Pydantic validation} for typed request/response schemas, improving robustness and reducing runtime errors.
  \item \textbf{Automatic OpenAPI documentation} via Swagger UI (\texttt{/docs}) and ReDoc (\texttt{/redoc}), supporting both assessment documentation and viva demonstrations.
\end{itemize}

\subsection{Database: PostgreSQL with pgvector}
A core design decision was to use PostgreSQL extended with \texttt{pgvector} rather than introducing a separate vector database. This provides:
\begin{itemize}[leftmargin=1.2em]
  \item \textbf{Unified datastore} for relational entities (repositories, questions, analysis metadata) and vector embeddings (code chunks).
  \item \textbf{Relational integrity} with foreign keys linking repositories, chunks, and questions, improving traceability and auditability.
  \item \textbf{Efficient similarity search} (cosine distance) directly within SQL for top-$k$ retrieval during RAG.
\end{itemize}
SQLAlchemy 2.0 (async) was used as the ORM, with Alembic for migrations.

\subsection{AI Models: Local Embeddings and Optional Generation}
To avoid paid APIs by default and support offline usage:
\begin{itemize}[leftmargin=1.2em]
  \item \textbf{Embeddings:} HuggingFace sentence-transformers \texttt{all-MiniLM-L6-v2} produces 384-dimensional embeddings for code chunks.
  \item \textbf{Optional LLM:} Ollama can be enabled to run a small local model (\texttt{mistral}) for grounded answer generation.
\end{itemize}
The system is designed so that Lightweight Mode does not require the LLM. When the LLM is unavailable, the Q\&A endpoint returns a clear service-unavailable outcome but still provides retrieved sources as evidence.

\subsection{Containerisation and Lightweight Mode}
The application is managed via Docker Compose. By default, only \texttt{api} and \texttt{db} start in Lightweight Mode. Optional profiles enable advanced components:
\begin{itemize}[leftmargin=1.2em]
  \item \textbf{\texttt{--profile async}} enables Celery + Redis for background ingestion.
  \item \textbf{\texttt{--profile llm}} enables Ollama for local generation.
\end{itemize}
This ensures a deterministic, storage-safe default while allowing richer functionality when hardware permits.

\section{System Architecture}

\subsection{Overview}
CodeQuery uses a layered architecture:
\begin{itemize}[leftmargin=1.2em]
  \item \textbf{API layer (FastAPI):} routes for repositories, questions, and analysis.
  \item \textbf{Service layer:} ingestion, chunking, embedding, and Q\&A engine.
  \item \textbf{Persistence layer (PostgreSQL + pgvector):} stores repositories, code chunks (with vectors), questions, and analysis state.
\end{itemize}

\subsection{Key Resources and Endpoints}
The API follows resource-based REST design. Table~\ref{tab:endpoints} summarises core endpoints.

\begin{table}[h]
\centering
\small
\begin{tabular}{@{}llp{9.2cm}@{}}
\toprule
\textbf{Method} & \textbf{Endpoint} & \textbf{Description} \\
\midrule
POST & \texttt{/api/v1/repositories} & Add repository for analysis (triggers ingestion). \\
GET & \texttt{/api/v1/repositories} & List repositories. \\
GET & \texttt{/api/v1/repositories/\{id\}} & Get repository details. \\
PUT & \texttt{/api/v1/repositories/\{id\}} & Update / re-analyse repository. \\
DELETE & \texttt{/api/v1/repositories/\{id\}} & Delete repository and related artefacts. \\
POST & \texttt{/api/v1/repositories/\{id\}/questions} & Ask a question about a repository. \\
GET & \texttt{/api/v1/repositories/\{id\}/questions} & List question history for a repository. \\
GET & \texttt{/api/v1/questions/\{id\}} & Retrieve a single question record. \\
DELETE & \texttt{/api/v1/questions/\{id\}} & Delete a question record. \\
GET & \texttt{/api/v1/repositories/\{id\}/analysis} & View analysis / ingestion status. \\
GET & \texttt{/api/v1/repositories/\{id\}/statistics} & Retrieve code statistics for a repository. \\
GET & \texttt{/api/v1/health} & Health check. \\
\bottomrule
\end{tabular}
\caption{Core CodeQuery API endpoints (as documented in the project README).}
\label{tab:endpoints}
\end{table}

\section{Repository Ingestion Pipeline}

\subsection{Clone and File Discovery}
When a repository is added, the system clones it and discovers files while filtering irrelevant directories (e.g., \texttt{node\_modules}, \texttt{.git}) to reduce noise. Lightweight Mode includes a storage warning and recommends testing with small repositories (e.g., \texttt{octocat/Hello-World}) to prevent excessive disk usage.

\subsection{Chunking}
The ChunkingService aims to preserve semantic units where possible. Python files are chunked using AST parsing to keep functions/classes intact, while other languages fall back to generic chunking. Each chunk stores metadata (file path and line spans) to support traceable citations.

\subsection{Embedding and Storage}
Each chunk is embedded using \texttt{all-MiniLM-L6-v2} and stored in PostgreSQL with its vector representation. This enables semantic search within SQL using pgvector similarity queries.

\section{Repository Q\&A and RAG Workflow}

\subsection{Retrieval}
When a question is posted, the system embeds the query and retrieves the top-$k$ most similar code chunks from pgvector. Retrieval is scoped to the target repository, ensuring evidence comes from the correct codebase.

\subsection{Augmented Generation}
If the LLM profile is enabled, a strict prompt containing the retrieved chunks and the user query is sent to the local model via Ollama to generate a grounded answer. The response is returned alongside sources that cite file paths and line spans.

\subsection{Graceful Fallback in Lightweight Mode}
If the LLM is unavailable (Lightweight Mode), the Q\&A engine does not crash. Instead, it returns a clear outcome indicating generation is unavailable (HTTP 503) while still returning the retrieved evidence. This preserves usefulness as a semantic code-search tool while keeping the system storage-safe by default.

\subsection{Persistence and Auditability}
All Q\&A interactions are persisted with answer text, confidence score, model used, processing time, and retrieved sources. This supports auditability (what was answered and based on which evidence) and allows users to review question history.

\section{Testing Strategy}
Automated testing was implemented using \texttt{pytest}. The test suite covers:
\begin{itemize}[leftmargin=1.2em]
  \item \textbf{API tests:} health checks, repository CRUD behaviours, correct handling of validation errors (422) and missing resources (404).
  \item \textbf{Service tests:} chunking behaviour and ingestion file discovery rules.
  \item \textbf{Reproducibility:} tests can be executed inside the container (\texttt{docker compose exec api pytest -q}) for examiner-friendly verification.
\end{itemize}
Testing improves reliability and demonstrates professional practice aligned with REST API development standards.

\section{Challenges and Lessons Learned}
\begin{itemize}[leftmargin=1.2em]
  \item \textbf{Balancing ambition with deployability:} A full RAG stack can become resource-heavy. Implementing Lightweight Mode with optional Docker profiles was critical to keep the default footprint small and reliable.
  \item \textbf{Async database management:} Async SQLAlchemy improves throughput but requires careful session lifecycle handling to avoid connection leaks and inconsistent transaction states.
  \item \textbf{Chunking trade-offs:} Larger chunks preserve context but reduce retrieval precision, while smaller chunks increase precision but can lose meaning. AST-based chunking for Python demonstrated the value of language-aware segmentation.
\end{itemize}

\section{Limitations and Future Work}
\subsection{Limitations}
\begin{itemize}[leftmargin=1.2em]
  \item \textbf{Repository size constraints:} Large repositories can increase ingestion time and disk usage on consumer hardware.
  \item \textbf{Local model performance:} CPU-only local generation is slower than hosted inference; quality is model-dependent.
  \item \textbf{Language-aware chunking coverage:} AST-based chunking is strongest for Python; other languages currently rely more on generic strategies.
\end{itemize}

\subsection{Future Work}
\begin{itemize}[leftmargin=1.2em]
  \item Add API key authentication for secure public deployment.
  \item Extend language-aware parsing/chunking to TypeScript/Java/Go.
  \item Improve progress reporting and cancellation for long analyses.
  \item Promote async ingestion as default for heavy repositories (Celery/Redis), keeping Lightweight Mode as a fallback.
\end{itemize}

\section{Generative AI Declaration}
In accordance with University guidance, generative AI tools (specifically Google Gemini / Antigravity Agent) were used for creative exploration and advanced problem-solving, aligning with the highest-level usage band:
\begin{itemize}[leftmargin=1.2em]
  \item \textbf{Concept Exploration \& Alternative Architectures:} GenAI was used extensively to evaluate and pivot between different design choices, such as contrasting a pure FastAPI setup versus incorporating Celery for async ingestion, and designing the innovative ``Lightweight Mode'' fallback to ensure the system is deployable on constrained hardware without crashing.
  \item \textbf{Data Structure \& RAG Prototyping:} AI tools assisted in designing the PostgreSQL \texttt{pgvector} schema and prototyping the semantic retrieval pipeline, allowing me to focus on curating the high-level logic and cutting-edge database architecture.
  \item \textbf{Responsibility:} While AI guided the creative conception of cutting-edge architectural patterns, all final design decisions, API routing, data modelling, and the \texttt{pytest} validation suite were verified and synthesised independently to meet coursework objectives.
\end{itemize}

These advanced architectural explorations firmly ground the project in modern, state-of-the-art backend practices. Thoughtful analysis of this usage is evident in the included logs. Examples of exported conversation logs demonstrating this creative, high-level structural usage are attached as supplementary material.
\end{document}